\def\stat{sind+bolotova}



\def\tit{ИСПОЛЬЗОВАНИЕ ПАДЕЖНОЙ ГРАММАТИКИ ПРИ~ИНФОРМАЦИОННОМ ПОИСКЕ\\

В~БАЗЕ ЗНАНИЙ ЭКСПЕРТНОЙ СИСТЕМЫ\\ О~КОНСТРУКЦИЯХ ЛЕТАТЕЛЬНЫХ 

АППАРАТОВ}



\def\titkol{Использование падежной грамматики при~информационном поиске 

в~базе знаний} % экспертной системы о~конструкциях летательных аппаратов}



\def\aut{Н.\,И.~Сидняев$^1$, Ю.\,И.~Бутенко$^2$, Е.\,Е.~Синева$^3$}



\def\autkol{Н.\,И.~Сидняев, Ю.\,И.~Бутенко, Е.\,Е.~Синева}



\titel{\tit}{\aut}{\autkol}{\titkol}



\index{Сидняев Н.\,И.}

\index{Бутенко Ю.\,И.}

\index{Синева Е.\,Е.}

\index{Sidnyaev N.\,I.}

\index{Butenko Yu.\,I.}

\index{Sineva E.\,E.}



%{\renewcommand{\thefootnote}{\fnsymbol{footnote}}

%\footnotetext[1]

%{Работа выполнена при частичной финансовой поддержке РФФИ 

%(проект 20--07--00655);

%исследования проводились в~рамках программы Московского центра фундаментальной 

%и~прикладной математики.}} 



\renewcommand{\thefootnote}{\arabic{footnote}}

\footnotetext[1]{Московский государственный технический университет им.\ Н.\,Э.~Баумана,\\ 

\mbox{sidnyaev@yandex.ru}}

\footnotetext[2]{Московский государственный технический университет им.\ Н.\,Э.~Баумана,\\ 

\mbox{iuliiabutenko2015@yandex.ru}}

\footnotetext[3]{Московский государственный технический университет им.\ Н.\,Э.~Баумана,\\

 \mbox{bolotovaee@mail.ru}}



 \label{st\stat}



 \thispagestyle{headings}

 

 \vspace*{-12pt} 

 

  

  

  

  \Abst{Раскрыты основные принципы построения экспертных систем в~области 

авиационной промышленности. Предложен ло\-ги\-ко-линг\-ви\-сти\-че\-ский подход 

к~информационному моделированию предметной области при проектировании базы знаний 

экспертной системы. Разработана модель представления знаний для описания конструкций 

летательных аппаратов, обеспечивающая информационно-структурную надежность базы 

знаний экспертной системы. Проведен анализ парадигматических особенностей 

авиакосмической терминологии с~целью улучшения эффективности информационного 

поиска. Предложен метод разрешения лексической многозначности поискового запроса 

в~базе знаний. Описан механизм вывода информации из базы знаний на основе падежной 

грамматики Ч.~Филлмора, позволяющей определить се\-ман\-ти\-ко-син\-таксическую структуру 

выводимого предложения. Показано, что посредством расстановки ограничений на число 

участников ситуации, которые может присоединять глагол, можно обеспечить эффективный 

вывод информации из базы знаний.}

  

  \KW{экспертная система; база знаний; летательный аппарат; терминология; знания}

  

\DOI{10.14357\!/\!08696527210307} 



\vskip 10pt plus 9pt minus 6pt



 \thispagestyle{headings}

 

 \vspace*{-16pt}

  

     \section{Введение}

     

     \vspace*{-2pt}

     

  Поддержка принятия решений в~сложных ситуациях необходима 

в~различных областях деятельности, где требуются обоснованные, логически 

доказуемые аргументы, в~том числе в~оперативном и~стратегическом 

управлении технологиями на авиационных предприятиях. При принятии 

решений необходимо проведение детального моделирования последствий 

предполагаемого решения, поиск оптимального пути достижения заданного 

результата при помощи технологий имитационного моделирования. 

С~подобными задачами успешно справляются экспертные системы~--- 

комплексы программных средств, способные заменять специалиста при 

решении сложных задач, возникающих в~процессе диагностики, 

проектирования или эксплуатации аэрокосмической и~другой техники~[1]. 

Такие системы основаны на знаниях, полученных в~процессе взаимодействия 

с~экспертами в~конкретной предметной об\-ласти.

  

  Параллельно с~развитием различных направлений исследований в~области 

искусственного интеллекта идет развитие информационных структур для 

представления знаний~[2, 3]. Появились новые способы описания 

и~представления данных, возникли фреймовые, списочные и~иерархические 

структуры. Процесс построения баз знаний достаточно сложен и~носит 

итеративный характер, заключающийся в~циклической модификации баз 

знаний на основе результатов ее тестирования~[4]. 



\vspace*{-6pt}



\section{Абстрактная иерархическая структура как способ описания 

конструкций летательных аппаратов}



\vspace*{-2pt}

  

  Основной задачей проектирования базы знаний в~области авиакосмического 

приборостроения является описание структуры летательного аппарата таким 

образом, чтобы оно предоставляло наиболее полную и~непротиворечивую 

информацию об описываемом объекте. Если применить лингвистические 

модели в~совокупности с~математическим аппаратом, можно достичь более 

ясного и~результативного описания и~решения сложной технологической 

проблемы. Информация воспринимается специалистом на когнитивном уровне 

посредством основных единиц этого уровня~--- понятий, идей, концептов, 

поэтому для представления знаний часто применяется аппарат семантических 

сетей, который позволяет воссоздать модель, схожую с~процессами мышления 

конструктора, технолога~[5, 6]. В инженерии знаний под семантической сетью 

подразумевается граф с~узлами, отображающими объекты предметной области, и~дугами,

 обозначающими отношения между данными понятиями. 

  

  Для описания конструкций летательных аппаратов разработана 

интегрированная модель представления знаний~--- абстрактная иерархическая 

структура, объединяющая графовую структуру, фреймы и~предикаты. Все 

знания о предметной области представлены в~виде графа, в~узлах которого 

содержатся фреймы~--- минимально возможные описания сущностей, 

записанные с~помощью предикатов и~предикатных символов~[1]. Подобная 

структура гарантирует непротиворечивость описаний объектов и~их отношений в~реальном производстве.



\vspace*{-6pt}



\section{Отражение парадигматических особенностей авиакосмической 

терминологии в~базе знаний}



\vspace*{-2pt}



  Неотъемлемой частью любого управления является его документационное 

обеспечение, основа которого~--- однозначно понимаемая и~непротиворечивая 

терминология. Единое терминологическое обеспечение позволяет избегать 

недопонимания между специалистами как в~рамках одной предметной области, 

так и~за ее пределами.

  

  Терминосистема~--- это некоторая лингвистическая модель, представляющая 

конкретную предметную область~[7]. Аэрокосмическая терминология русского 

языка представляет собой сложную систему, включающую термины из 

различных тематических пространств: обозначения субъектов аэрокосмической 

сферы (\textit{космонавт}, \textit{летчик}, \textit{пилот}), наименования 

организаций и~структур области авиации и~космонавтики (\textit{НАСА}, 

\textit{ИАТА}, \textit{ИФАТКА}), обозначения действий аэрокосмической 

деятельности (\textit{взлет}, \textit{посадка}, \textit{пилотирование}, 

\textit{пристыковка}), наименования деталей и~частей летательных аппаратов 

(\textit{фюзеляж}, \textit{крыло}, \textit{подкрылки}, \textit{двигатель}), 

обозначение характеристик и~состояний летательных аппаратов 

(\textit{работоспособность}, \textit{надежность}, \textit{безопас\-ность}). 

Стандартизация терминологии данной предметной области~--- процесс 

трудоемкий, а~отсутствие стандартизации приводит к~появлению разных 

терминов для обозначения одних и~тех же объектов и~явлений, в~связи с~чем 

возникают случаи синонимии и~лексической многозначности.

  

  При построении баз знаний интеллектуальных систем синонимия считается 

крайне неблагоприятным явлением. Различия в~номинации объектов приводят 

к~трудностям в~сфере профессиональной коммуникации, поэтому базу знаний 

экспертной системы предлагается пополнять всеми возможными вариантами 

синонимов тех или иных узкоспециальных и~общенаучных терминов.

  

  В качестве иллюстрации приведем несколько примеров синонимичных 

терминов: \textit{ин\-тер\-цеп\-тор}--\textit{спой\-лер}; \textit{пусковое 

устрой\-ст\-во}\,--\,\textit{пусковая кнопка}\,--\,\textit{стар\-тер}; \textit{шкала 

даль\-ности}\,--\,\textit{дис\-тан\-ци\-он\-ная шкала}; \textit{спуск}--\textit{сни\-же\-ние}; 

\textit{приборная дос\-ка}\,--\,\textit{при\-бор\-ная панель}; 

\textit{мас\-ло\-уло\-ви\-тель}--\textit{мас\-ло\-сбор\-ник}.

  

  В аэрокосмической терминологии также наблюдается семантическая 

оппозиция номинаций <<свой--чу\-жой>>, в~результате чего возникают 

синонимичные пары, например: \textit{кос\-мо\-навт}--\textit{аст\-ро\-навт} 

(\textit{англ}.\ astrounaut); \textit{кос\-мо\-нав\-ти\-ка}--\textit{аст\-ро\-нав\-ти\-ка} 

(\textit{англ}.\ astrounautica); \textit{вы\-со\-то\-мер}--\textit{аль\-ти\-метр} (\textit{франц}.\ 

altim\`{e}tre).

  

  Равнозначность подобных синонимов необходимо отразить в~проектируемой 

базе знаний во избежание недопонимания между экспертной сис\-те\-мой и~ее 

пользователями. Стоит отметить, что редко используемые аббревиатуры 

должны быть расшифрованы и~понятны пользователям экспертной сис\-те\-мы 

(например, \textit{ШВП}~--- это \textit{шасси на воздушной подушке} или 

\textit{шаровой вытяжной парашют}).

  

  Значительные трудности в~процессе поиска информации в~базе знаний 

вызывают явления лексической многозначности, т.\,е.\ сосуществование 

различных значений одного слова или фразы, например:\\[-16pt]



\begin{description}

\item[\,]

\textit{Свеча}~--- 1.~Приспособление для воспламенения горючей смеси 

(в~двигателях внутреннего сгорания); 2.~Крутой взлет, подъем (фигура 

высшего пилотажа).\\[-14.5pt]



\item[\,]

\textit{Хвост}~--- 1.~Хвостовое оперение летательных аппаратов; 2.~Вытянутое 

из пыли и~газа кометное вещество, образующееся при приближении кометы 

к~Солнцу.\\[-14.5pt]



\item[\,]

\textit{Атмосфера}~--- 1.~Газообразная оболочка, окружающая Землю 

и~некоторые другие планеты; 2.~Единица измерения дав\-ле\-ния.

\end{description}



\pagebreak



\section{Применение падежной грамматики Ч.~Филлмора для вывода 

знаний из экспертной системы}



\vspace*{-2pt}

   

   Один из необходимых компонентов базы знаний экспертной системы~--- 

механизм вывода, который содержит правила для решения конкретной задачи. 

Механизм вывода ссылается на информацию из базы знаний и~выбирает факты 

и~правила, которые будут применяться при попытке ответить на запрос 

пользователя. Данный механизм обеспечивает аргументацию информации 

в~базе знаний и~позволяет человеку получать информацию из базы знаний 

в~виде предложений, построенных на естественном языке. Одним из наиболее 

удачных механизмов, позволяющих моделировать смысл высказываний на 

естественном языке, считается механизм, основанный на теории падежных 

грамматик Ч.~Филлмора~[8]. В~падежной грамматике глагол является центром 

предложения и~диктует в~силу своего значения набор ролей (глубинных 

падежей), реализующихся в~предложении посредством именных форм. Иными 

словами, падежная грамматика показывает связь существительного или 

местоимения с~другими словами в~предложении. Обрамление падежами 

подвержено определенным ограничениям, например каж\-дый семантический 

падеж может встречаться в~предложении только один раз~[9]. Некоторые 

падежи являются обязательными, другие~--- необязательными. Обязательные 

падежи нельзя удалять, рискуя получить грамматически неправильные 

предложения.

   

   Для описания события в~каждом фрейме иерархической структуры в~первую 

очередь выделяется действие, которое обычно описывается глаголом. Далее 

определяются: объект, который действует~--- агенс; объект, над которым это 

действие выполняется~--- пациенс~[10]. 

  

  На каждое отношение накладывается множество ограничений, например 

структура глагола <<пилотировать>> может включать следующие 

семантические падежи: (1)~агентив~--- {\bfseries\textit{летчик}} 

\textit{пилотировал самолет}; (2)~пациентив~--- \textit{летчик пилотировал} 

{\bfseries\textit{самолет}}; (3)~инструменталис~--- \textit{летчик пилотировал 

самолет} {\bfseries\textit{боковой ручкой управления}}; (4)~локатив~--- 

\textit{летчик впервые пилотировал самолет} {\bfseries\textit{на аэродроме}} 

\textit{с~плохим покрытием}. Возможно включение в~структуру данного 

глагола и~других семантических падежей, а~также атрибутивных отношений, 

например: <<\textit{пилотировать каким образом}?>>~--- <<\textit{умело}>>, 

<<\textit{профессионально}>>.

  

  Такие ограничения необходимо накладывать для того, чтобы сис\-те\-ма могла 

стро\-ить правильные се\-ман\-ти\-ко-син\-так\-си\-че\-ские структуры, что 

обеспечит связ\-ность и~внут\-рен\-нюю ин\-тер\-пре\-ти\-ру\-емость знаний в~экспертной 

сис\-теме.

  

  Падежные отношения, а также соответствие между синтагматикой 

и~линейной упорядоченностью слов в~предложении можно представить в~виде 

дерева непосредственно составляющих. Глубинная структура 

пропозиционального компонента любого простого предложения представляет 

собой конструкцию, состоящую из глагольной и~именной группы, которые 

находятся в~специальных помеченных отношениях (семантических падежах) ко 

всему предложению. Дерево составляющих позволяет представить 

синтаксическую структуру предложения в~виде иерархии, составленной из 

связанных друг с~другом не\-пе\-ре\-се\-ка\-ющих\-ся единиц. Пример описания 

предложения <<\textit{Грузовой космический корабль произвел стыковку с~МКС 

в~штатном режиме}>> с~помощью дерева составляющих показан на рис.~1. 

  

  \begin{figure}  %fig1

  \vspace*{1pt}

\begin{center}

 \mbox{%

 \epsfxsize=97.104mm 

\epsfbox{bol-1.eps}

 }

\end{center}

\vspace*{-9pt}

  \Caption{Представление предложения в~виде дерева со\-став\-ля\-ющих}

  \vspace*{-6pt}

  \end{figure}

   

   Исходя из этого, модель семантической целостности предложения ($C$) 

можно представить сле\-ду\-ющим образом:

   $$

   C=\langle V^I, N^I\rangle\,,

   $$

где $V$~--- глагольная составляющая; $N$~--- именная составляющая; $I$~--- 

семантический падеж. Для вывода информации из базы знаний требуются 

также атрибутивные отношения для описания формы, цвета, размера и~других 

характеристик объектов, а~так\-же дополнительные падежи цели, условия, 

времени и~др., т.\,е.\ в~за\-ви\-си\-мости от целей экспертной сис\-те\-мы перечень, 

содержащий 6~основных падежей, может быть расширен.

   

   Добавим к~структуре рассматриваемого предложения дополнительные 

семантические падежи и~применим теорию семантических падежей для 

описания семантической структуры знания о~событии. Допустим, системе 

нужно вывести сообщение <<\textit{Грузовой космический корабль 

пристыковался к~МКС в~15:00~мск в~штатном режиме. Груз доставлен на 

станцию}>>. Семантическая сеть, опи\-сы\-ва\-ющая событие до\-став\-ки груза на 

МКС, пред\-став\-ле\-на на рис.~2.

     

     \begin{figure} %fig2

     \vspace*{1pt}

\begin{center}

 \mbox{%

 \epsfxsize=107.782mm 

\epsfbox{bol-2.eps}

 }

\end{center}

\vspace*{-9pt}

     \Caption{Пример семантической сети, опи\-сы\-ва\-ющей событие доставки груза на 

МКС}

\vspace*{-6pt}

      \end{figure}

  

  В процессе поиска информации в~базе знаний, основанной на семантической 

сети, происходит сопоставление общей сети с~сетью запроса для поиска 

нужного фрагмента информации. При выводе в~семантических сетях также 

применяется метод перекрестного поиска, в~процессе которого сопоставляются 

не узлы семантической сети, а~дуги~--- отношения между объектами.



\vspace*{-6pt}

     

     \section{Заключение}

     

     \vspace*{-2pt}

  

  В работе раскрыты основные принципы построения интеллектуальных 

сис\-тем в~авиакосмической промышленности. Поднимается проблема 

стандартизации терминологии как одного из ключевых этапов разработки баз 

знаний экспертных систем с~целью сведения разнообразно представленной 

информации к~общепонятному виду, что позволит организовать ее совместное 

и~многократное использование специалистами. Для выявления 

парадигматических особенностей терминологии авиакосмической области 

проведен анализ лексических единиц, входящих в~тематическую группу 

<<Авиакосмическое приборостроение>>. В~связи с~тем что стандартизация 

терминологии является трудоемким процессом, предложено пополнять 

проектируемую базу знаний всеми возможными вариантами синонимичных 

терминов. Рассмотрен механизм вывода информации из базы знаний 

экспертной системы на основе падежной грамматики Ч. Филлмора, которая 

позволяет определить се\-ман\-ти\-ко-син\-так\-си\-че\-скую структуру 

выводимого предложения. Показано, что посредством расстановки 

ограничений на число актантов (участников ситуации), которые может 

присоединять глагол, можно обеспечить эффективный вывод предложений на 

естественном языке. 



\vspace*{-6pt}



{\small\frenchspacing

{%\baselineskip=10.8pt

\addcontentsline{toc}{section}{Литература}

\begin{thebibliography}{99}



\vspace*{-2pt}



\bibitem{1-bol}

\Au{Сидняев~Н.\,И., Бутенко~Ю.\,И., Болотова~Е.\,Е.} Экспертная система продукционного 

типа для создания базы знаний о конструкциях летательных аппаратов~/\!\!/ Авиакосмическое 

приборостроение, 2019. №\,6. С.~38--52. doi: 10.25791\!/\linebreak aviakosmos.06.2019.676.



\bibitem{3-bol} %2

\Au{Сигов~А.\,С., Нечаев~В.\,В., Кошкарёв~М.\,И.} Архитектура пред\-мет\-но-ори\-ен\-ти\-ро\-ван\-ной 

базы знаний интеллектуальной сис\-те\-мы~/\!\!/ Int. J.~Open Information 

Technologies, 2014. No.\,12. P.~1--6.

\bibitem{2-bol} %3

\Au{Сидняев~Н.\,И., Бутенко Ю.\,И., Болотова~Е.\,Е.} Логическая модель требований 

ин\-фор\-ма\-ци\-он\-но-сис\-тем\-ной на\-деж\-ности для баз знаний интеллектуальных сис\-тем~/\!\!/ 

Программная инженерия, 2020. №\,4. С.~195--204. doi: 10.17587\!/\!prin.11.195-204.



\bibitem{4-bol}

\Au{Helbig~H.} Knowledge representation and the semantics of natural language.~--- 

Berlin/Heidelberg/New York: Springer, 2006. 655~p.

\bibitem{5-bol}

\Au{Woods~D.\,D.} Cognitive technologies: The design of joint human-machine cognitive systems~/\!\!/ 

AI Mag., 1985. Vol.~6. Iss.~4. P.~86--92.

\bibitem{6-bol}

\Au{Ломов~П.\,А., Шишаев~М.\,Г.} Формирование когнитивных фреймов на основе 

онтологических паттернов для визуализации онтологий~/\!\!/ Информационные сис\-те\-мы 

и~технологии, 2015. №\,6(92). С.~12--22. 

\bibitem{7-bol}

\Au{Даниленко~В.\,П.} Русская терминология: опыт лингвистического описания.~--- М.: Наука, 

1977. 246~с.

\bibitem{8-bol}

\Au{Fillmore~Ch.} The case for case~/\!\!/ Universals in linguistic theory~/

Eds. E.~Bach, R.\,T.~Harms.~---

 New York, NY, USA: Holt, 

Rinehart, and Winston, 1968. P.~1--88. 

\bibitem{9-bol}

\Au{Богданов~В.\,В.} Струк\-тур\-но-се\-ман\-ти\-че\-ская организация предложения.~---

 Л.: ЛГУ, 1977. 205~с.

\bibitem{10-bol}

\Au{Jackendoff~R.\,S.} The status of thematic relations in linguistic theory~/\!\!/ Linguistic 

Inq., 1987. Vol.~16. P.~369--411.

  \end{thebibliography}

} }







\vspace*{-8pt}



\hfill{\small\textit{Поступила в~редакцию 18.08.21}}



\vspace*{6pt}



\hrule



\vspace*{2pt}



%\newpage



%\vspace*{-28pt}



\hrule



%\vspace*{8pt}



\def\tit{USE OF CASE GRAMMAR IN~INFORMATION SEARCH IN~THE~EXPERT SYSTEM KNOWLEDGE 

BASE\\ ON~AIRCRAFT STRUCTURES}





\def\titkol{Use of case grammar in~information search in~the~expert system knowledge 

base} % on~aircraft structures}



\def\aut{N.\,I.~Sidnyaev, Yu.\,I.~Butenko, and E.\,E.~Sineva}



\def\autkol{N.\,I.~Sidnyaev, Yu.\,I.~Butenko, and E.\,E.~Sineva}



\titel{\tit}{\aut}{\autkol}{\titkol}



\vspace*{-11pt}



\noindent

Bauman Moscow State Technical University, 5-1~2nd Baumanskaya Str., Moscow 105005, Russian 

Federation



\def\leftfootline{\small{\textbf{\thepage}

\hfill Sistemy i Sredstva Informatiki~---

Systems and Means of Informatics 2021 vol~31 no\,3}

}%

 \def\rightfootline{\small{Sistemy i Sredstva Informatiki~---

 Systems and Means of Informatics 2021 vol~31 no\,3

\hfill \textbf{\thepage}}}



\vspace*{3pt}







\Abste{The paper reveals the basic principles of building expert systems in the aviation industry. 

A~model of knowledge representation that provides information and structural reliability of the expert 

system's knowledge base for describing aircraft structures has been developed. 

The article analyzes the paradigmatic features of aerospace terminology in order to 

improve the efficiency of information search. A~method for resolving lexical ambiguity of 

a~search query in the knowledge base is proposed. Inference engine and explanation mechanism

of 

the knowledge base premised on the Charles Fillmore's case grammar is described. 

The grammatical cases make it possible to determine the semantic and syntactic structure

 of the output sentence. Restrictions show the number of participants in the situation that 

 a~verb can attach, it is possible to 

provide effective output of information from the knowledge base.}



\KWE{expert system; knowledge base; aircraft; terminology; knowledge}



\DOI{10.14357\!/\!08696527210307} 



%\vspace*{-24pt}



%\Ack

%\noindent





\vspace*{-12pt}



\renewcommand{\bibname}{\protect\rmfamily References}

%\renewcommand{\bibname}{\large\protect\rm References}



{\small\frenchspacing

{%\baselineskip=10.8pt

\addcontentsline{toc}{section}{References}

\begin{thebibliography}{99}

\bibitem{1-sin-1}

\Aue{Sidnyaev, N.\,I., Yu.\,I.~Butenko, and E.\,E.~Bolotova.} 2019. Ekspertnaya sis\-te\-ma 

pro\-duk\-tsi\-on\-no\-go tipa 

dlya so\-zda\-niya bazy zna\-niy o~konstruktsiyakh le\-ta\-tel'\-nykh ap\-pa\-ra\-tov [Rule-based expert system for creating 

a~knowledge base on aircraft structures]. \textit{Aviakosmicheskoe priborostroenie} [Aerospace

 Instrument-Making] 6:38--52. doi: 10.25791\!/\!aviakosmos.06.2019.676.



\bibitem{3-sin-1}

\Aue{Sigov, A.\,S., V.\,V.~Nechaev, and M.\,I.~Koshkaryov.}

 2014. Arkhitektura predmetno-oriyentirovannoy 

ba\-zy znaniy in\-tel\-lek\-tu\-al'\-noy sis\-te\-my [Architecture of the subject-oriented knowledge base of the intellectual 

system]. \textit{Int. J.~Open Information Technologies} 12:1--6. 



\bibitem{2-sin-1}

\Aue{Sidnyaev, N.\,I., Yu.\,I.~Butenko, and E.\,E.~Bolotova.} 2020. Lo\-gi\-che\-skaya model' tre\-bo\-va\-niy 

informatsionno-sistemnoy na\-dezh\-nosti dlya baz zna\-niy in\-tel\-lek\-tu\-al'\-nykh sis\-tem [Logical model of 

information and system reliability requirements for knowledge bases of intellectual systems]. 

\textit{Programmnaya inzheneriya} [Software Engineering] 4:195--204. doi: 10.17587\!/\!prin.11.195-204.



\bibitem{4-sin-1}

\Aue{Helbig, H.} 2006. Knowledge representation and the semantics of natural language. 

Berlin/Heidelberg/New York: Springer. 655~p.

\bibitem{5-sin-1}

\Aue{Woods, D.\,D.} 1985. Cognitive technologies: The design of joint human-machine cognitive systems. 

\textit{AI Mag.} 6(4):86--92.

\bibitem{6-sin-1}

\Aue{Lomov, P.\,A., and M.\,G.~Shishaev.} 2015. Formirovanie kog\-ni\-tiv\-nykh frey\-mov na osno\-ve 

on\-to\-lo\-gi\-che\-skikh pat\-ter\-nov dlya vi\-zu\-a\-li\-za\-tsii on\-to\-lo\-giy 

[Approach to

ontology visualization based cognitive frames]. 

\textit{Informatsionnye sistemy i~tekhnologii} [Information  Systems and Technologies] 6(92):12--22.

\bibitem{7-sin-1}

\Aue{Danilenko, V.\,P.} 1977. \textit{Russkaya terminologiya: opyt lingvisticheskogo opisaniya} [Russian 

terminology: Experience of linguistic description]. Moscow: Nauka. 246~p.

\bibitem{8-sin-1}

\Aue{Fillmore, C.\,J.} 1968. The case for case. \textit{Universals in linguistic theory}. 

Eds. E.~Bach and R.\,T.~Harms. 

New York, NY: Holt, Rinehart, and Winston. 1--88. 

\bibitem{9-sin-1}

\Aue{Bogdanov, V.\,V.} 1977. \textit{Strukturno-semanticheskaya or\-ga\-ni\-za\-tsiya pred\-lo\-zhe\-niya} 

[Structural and semantic organization of a~sentence]. Leningrad: LSU. 205~p.

\bibitem{10-sin-1}

\Aue{Jackendoff, R.\,S.} 1987. The status of thematic relations in linguistic theory. 

\textit{Linguistic Inq.} 16:369--411.



\end{thebibliography}

} }



\vspace*{-6pt}



\hfill{\small\textit{Received August 18, 2021}} 



\vspace*{-12pt}  



\Contr



\vspace*{-4pt}



\noindent

\textbf{Sidnyaev Nikolay I.} (b.\ 1955)~--- Doctor of Science in technology, professor, Head of Department, 

Bauman Moscow State Technical University, 5-1 2nd Baumanskaya Str., Moscow 105005, Russian 

Federation; \mbox{sid-nyaev@yandex.ru}



\vspace*{3pt}



\noindent

\textbf{Butenko Yuliya I.} (b.\ 1987)~--- Candidate of Science (PhD) in technology, associate professor, 

Bauman Moscow State Technical University, 5-1~2nd Baumanskaya Str., Moscow 105005, Russian 

Federation; \mbox{iuliiabutenko2015@yandex.ru}



\vspace*{3pt}



\noindent

\textbf{Sineva Elizaveta E.} (b.\ 1998)~--- Master student, Faculty of Fundamental Sciences, Bauman 

Moscow State Technical University, 5-1 2nd Baumanskaya Str., Moscow 105005, Russian Federation; 

\mbox{bolotovaee@mail.ru}





\label{end\stat}



\renewcommand{\bibname}{\protect\rm Литература} 

